\documentclass[11pt]{article}
\usepackage[a4paper, top=18mm, bottom=18mm, left=20mm, right=20mm]{geometry}
\usepackage[T2A]{fontenc}
\usepackage[utf8]{inputenc}
\usepackage[ukrainian]{babel}
\usepackage{graphicx}
\usepackage[export]{adjustbox}
\usepackage{amsmath}
\usepackage{amssymb}
\graphicspath{{/app/Quiz/Scripts/2023/ProgAll/15BinTree/}{/app/Quiz/Scripts/2021/Mod2021_3/BinTree/}{/app/Quiz/Scripts/2020/Mod2020_4/BinTree/}}
\setlength{\parindent}{0pt}
\setlength{\parskip}{6pt}
\pagestyle{empty}
\begin{document}
%begin
{\large\textbf{Контрольна робота}}

\textbf{Варіант \No\,4}\hfill \textit{ПІБ:}\,\underline{\hspace{60mm}}
\vspace{4mm}

%beginitem
1. Що буде виведено за виконання фрагменту коду?

print('R', end=' ')
for d in range(3,3,-3):
    if d>=6:
        break
        print(d, end=' ')
    if d>=6:
        break
    else:
        print(d, end=' ')
    print(d)
\vspace{5mm}
%enditem

%beginitem
2. Що буде виведено за виконання фрагменту коду?
a = 3
b = 1
c = 5

def h(a):
global c    a = 3
    b *= 2
    c = 4
    return a + b + c

a = 0
b = 9
c = 6

try:
    print(h(b), a, b, c, sep=';')
except:
    print('Z', a, b, c, sep=';')
\vspace{5mm}
%enditem

%beginitem
3. 

\begin{center}
	\adjustimage{max size={0.8\linewidth}{0.8\paperheight}}
	{../15BinTree/trees_exp/59.png}
\end{center}
Указати послідовність міток вершин дерева, зображеного на рис., що утвориться в результаті обходу дерева в прямому порядку. Мітки записувати через пробіл.\par Алгоритм починає роботу з кореня (на рис. це верхня вершина).
%answer:59 прямому
\vspace{5mm}
%enditem

%beginitem
4. 
Записати ВИРАЗ, значенням якого є список, що складається з кубів цілих чисел від 12 до 2018 (включно), що не діляться на жодне з чисел 2, 3, записаних в оберненому порядку.\vspace{5mm}
%enditem

%beginitem
5. %goodpath:Scripts/2019/Mod2019_1/Lex/01-good.txt
Відмітити все, що є однією правильною лексемою:
( 1 ) (1)

( 2 ) 0.2E3

( 3 ) "abc'

( 4 ) x23
\vspace{5mm}
%enditem

%beginitem
6. 
Написати програму, що запитує у користувача значення дійсних параметрів $x$ та $y$, обчислює для них значення виразу 
$$\frac{\log_{2} x + 12}{(\tg x + 0.3) (y + 68)}$$ 
та повiдомляє введенi данi та результат обчислення.
 Оцінюється правильність та якість коду.

\vspace{5mm}
%enditem

%beginitem
7. Що буде виведено за виконання фрагменту коду?
d=[10,11,12,13,14]; print(d.pop(-2),end=' '); print(d)\vspace{5mm}
%enditem

%beginitem
8. 
Записати ВИРАЗ, значенням якого є множина, що складається з цілих чисел від 12 до 2018 (включно), що не діляться на 7.\vspace{5mm}
%enditem

%beginitem
9. %goodpath:Scripts/2018/Mod2018_1/01-good.txt
Відмітити все, що є однією правильною лексемою:
( 1 ) 0,5

( 2 ) LIFO

( 3 ) x+23

( 4 ) x23
\vspace{5mm}
%enditem

%beginitem
10. %goodpath:Scripts/2019/Mod2019_1/Lex/03-good.txt
Відмітити все, що є коректним оператором С++:
( 1 ) if a<b: a=0 else: b=0

( 2 ) x=3**1

( 3 ) if ('x'><'y'): a=6

( 4 ) --n
\vspace{5mm}
%enditem

%end
\newpage
\end{document}

